\section{Proof of the equidistribution theorem}\label{sec:assembly}

We now return to the packets of Theorem~\ref{thm:main}.
Proposition~\ref{prop:cusp} and Mahler's compactness criterion give
tightness. Take an arbitrary subsequence with a weak probability limit
$\mu$. It is $A$-invariant, and Proposition~\ref{prop:ball} gives
the ordinary-ball estimate required by
Proposition~\ref{prop:positive-entropy}. Theorem~\ref{thm:ekl}
therefore makes almost every $A$-ergodic component of $\mu$
homogeneous.

Apply the intermediate-conditioning construction to this subsequence,
passing further to a subsequence wherever necessary. The augmented
limiting state still projects to $\mu$ on $X$.
Proposition~\ref{prop:bowen} supplies the tube estimate throughout
the required time interval. Hence
Theorem~\ref{thm:full-activity} makes every root active almost surely
in this state.

If proper homogeneous components had positive total weight,
Lemma~\ref{lem:fixed-pieces} would place positive $\mu$-mass on the
countable union of their fixed-point pieces. At least one piece would
have positive mass. Its fixing diagonal is nonidentity, so it has
two unequal diagonal entries. The corresponding root plaques meet
that piece in at most countably many points. By
Lemma~\ref{lem:transverse-inactive}, that root is inactive almost
surely on states projecting to the piece, a contradiction.
Thus almost every component is the $G$-invariant probability, and
$\mu=m_6$.

Every weakly convergent subsequence supplied by tightness has the same
limit. It follows that the original sequence converges to $m_6$ on
$C_c(X_6)$. Proposition~\ref{prop:cusp} already gives the compact sets
required for the nonescape assertion. The estimates are uniform in the
embedding ordering, so the conclusion holds for the arbitrary choices
in Theorem~\ref{thm:main}.
